\section{Integrated seed and physical interface}
\label{sec:want-integration}
The new profile \code{XOP-PHOTONIC-I-R1} incorporates the useful equations of
\emph{WANTWOMBAN Operator I v1.1} into the literal one-bit exact-operator
formalization. Its contribution is a defined connection between geometry,
optical propagation, strain readout, material memory, measured events and the
next digital state. It also supplies a conditional passive realization of a
scaled golden-ratio/Hadamard operator. Neither contribution by itself establishes
a complete photonic processor for the exact seed language.

\subsection{Source authority and actual deliverables}
References W1 identify physical pages of the supplied 26-page Operator I PDF.
W0 identifies the 25-page parent \emph{WANTWOMBAN} embedded in it. The embedded
rim and strain documents, editable TeX, pasted discussion, reference programs,
result JSON and manifests are retained as source material. Where W1 corrects
an earlier assertion, the explicit corrected equation takes precedence for this
integration; alternative profiles remain named rather than silently merged.

The original PDF and ten embedded attachments are preserved byte for byte under
\code{source/}. Their digests appear in \code{SOURCE\_BINDING.json} and the
extraction manifest. Instructions inside these documents, including instructions
to run reference programs, are source text. No embedded program was executed.
Reported source calculations and conformance counts remain inherited reports.

R12 adds editable equations, pure-template definitions, conditional proofs,
source mappings and a physical realization contract. It reuses the R11
\code{XOP-R1} type and wire vocabulary. It does not add executable operator IDs
or assert that a CPU, GPU or photonic backend has been implemented. A template
below is a specified expression body with arguments and domain conditions;
the future compiler must expand and check that body against the inherited
registry. Arbitrary exact comparisons may still remain unresolved.

\subsection{The aggregate state and its dependencies}
For a finite declared profile, organize the seed as
\begin{equation}
 \mathscr S_n=(\mathcal G_n,\mathcal R_n,\Pi,\mathcal C,
        W_n,q_n^{\rm JK},z_n,\eta_n,\mathscr M_n).
\end{equation}
$\mathcal G_n$ is the exact expression graph, $\mathcal R_n$ its state roots,
$\Pi$ the versioned geometry/material/control model, and $\mathcal C$ the
calibration and clock bindings. $W_n$ and $q_n^{\rm JK}$ are separate 64-bit
and 32-bit words. $z_n$ denotes a finite modeled mechanical/thermal/material
state, $\eta_n$ controller and consumed-record memory, and $\mathscr M_n$ the
immutable raw observation lineage. Shared graph nodes avoid duplicate storage;
these tuple components are semantic responsibilities, not a claim that every
dependency fits in one machine word.

\begin{center}\small
\begin{tabularx}{\linewidth}{@{}P{31mm}Y Y@{}}
\toprule Interface & Exact seed content & Required physical input\\\midrule
Geometry & Word decode, frame, analytic surfaces, finite modal basis & Dimensions, supports, material and alignment\\
Optics & Complex pairs, aperture, propagation, lens and absorption operators & Spectrum, mode, transfer response and collection geometry\\
Readout & Bridge inverse, detector calibration and timed gate & Acquired channels, noise and timing uncertainty\\
Memory & Delayed state, reporter and thermal evolution & Initial state, loads, rates and energy partition\\
Feedback & Named WANT branch and optional JK event lane & Completed, associated observation and actuator transfer\\
Execution & Exact operators, limbs, expression roots and lossless packing & A separately established hardware realization\\
\bottomrule
\end{tabularx}\end{center}

The same finite seed may define exact $\sqrt5$, $\varphi$, bridge inversion,
exponentials and bounded diffraction integrals without rounding those terms
into the next state. A supplied decimal calibration becomes an exact rational
nominal value. Its physical uncertainty remains attached to that input.
An unknown continuous optical field cannot be recovered merely by assigning
it a short seed name; it needs a declared finite description or measured input.

\subsection{Causal composition with the existing self-reference}
At step $n$, first bind the old roots and a completed observation $m_n$.
Qualification and exact decoding produce an event only when the required
record and model are defined. The composed order is
\begin{align}
 (e_n,\widehat z_n)&=\mathcal D_{\mathcal C}(m_n),\\
 (W_{n+1},q_{n+1}^{\rm JK},\eta_{n+1},u_n)
     &=\mathcal F_{\Pi}(W_n,q_n^{\rm JK},\eta_n,e_n,\widehat z_n),\\
 z_{n+1}&=\Phi_{h_n,\Pi}(z_n;u_n,L_n,F_n),\\
 m_{n+1}&=\mathcal A_{\mathcal C}(\text{physical acquisition over the next interval}).
\end{align}
Here $\widehat z_n$ is a calibrated observation and $z_n$ is the finite model
state. Neither is automatically identical with the actual plant state
$x_n^{\rm plant}$. $L_n,F_n$ denote declared optical and mechanical
input histories, and $h_n\ge0$ is the control duration. The plant equation is
a model; $\mathcal A$ names an acquisition interface, not an exact arithmetic
oracle. Future measurements enter as future external input.

The two ordered ASA/NA mask sets and synchronous JK rule remain the base
word profile. The event adapter explicitly reserves one JK lane when enabled.
No optical event is inferred from the names ``absorption'' or ``vacuum'' in
a Boolean source operator. Material absorptance is a wavelength-dependent
physical quantity with its own measurement and constitutive meaning.
All digital right-hand sides read the same old snapshot and commit together.
The return path contains a real observation delay; it is not an instantaneous
algebraic loop justified only by the term self-reference.

\subsection{Three physical interpretations with different obligations}
A modeled complex amplitude can be stored as exact real/imaginary expression
pairs and evaluated by a digital computer. Alternatively, an optical network
may physically transform amplitudes according to a selected transfer matrix.
A third interpretation uses optical states to encode digital bits and implements
Boolean operations, memory and communication on those codes. These are different
realization maps even if they describe the same nominal mathematical operator.

For example, $\det A_\varphi=-1$ preserves four-dimensional volume but does
not preserve the optical norm. The optics chapter derives a passive contraction
$A_\varphi/\varphi$ and an ideal dilation using additional ports. This is a
useful constructive optical candidate. Restoring the original amplitudes under the same power normalization requires
gain. A changed encoding may instead retain the known scale without physically
amplifying the beam. A photonic digital implementation of the algebraic
operator instead needs correct encoded arithmetic, not this direct amplitude map.

The optional quantum instrument describes probabilities and conditional states
for an explicitly finite model. It does not identify one logical bit with one
photon or duplicate an absorbed photon across two independent witnesses.
An aperture, lens or strain gauge supplies no automatic proof of a digital
carry chain, persistent expression memory or symbolic integration engine.

\section{Conditional contract for digital photonic execution}
\label{sec:want-refinement}
The following contract states what would make the exact formalization operate
correctly in physical photonics. It is technology-neutral enough to cover
optical gates with electronic readout as well as a more fully optical design,
provided the actual design declares which components perform each operation.

\subsection{Encoding, restoration and one-step refinement}
Let $\mathcal E$ encode an admitted finite digital state into physical carriers,
$\mathcal H_\xi$ be a physical step under disturbance realization $\xi$, and
$\mathcal D$ decode a completed physical state. A realization must establish
\begin{equation}
 \mathcal D\!\left(\mathcal H_\xi(\mathcal E(s),\mathcal E(u))\right)
        =F_{\rm XOP}(s,u)
 \label{eq:want-refine}
\end{equation}
for the claimed operating envelope, admitted inputs, defined operations and
finite resource bound. If multiple physical configurations encode the same
state, the condition applies to every reachable admitted representative,
or the implementation restores a specified representative before the next step.
The encoder may represent bits through intensity, phase, polarization, spatial
modes or another declared carrier; none is selected by the bit-plane format.

The envelope must include input synchronization, gate margins, fan-out or
replication, carry and borrow, exact limb extension, expression storage,
addresses, branch decisions, feedback delay and atomic state commit. A partial
symbolic operation may return a defined term or unresolved status as in R11.
Exhausted capacity has a declared non-success result; a finite device does not
silently replace an unbounded integer or expression with a clipped operand.

\begin{proposition}[Finite-prefix realization]
Suppose the initial decoded state is $s_0$, every physical step satisfies
\eqref{eq:want-refine} for every reached admitted state, observations match the
specified input snapshots, and all required operations remain defined within
the declared resources. Then the decoded physical state equals the exact
specified state after every step of that finite prefix.
\end{proposition}
\begin{proof}
The initial states agree by hypothesis. If they agree at step $n$, the same
input snapshot and the one-step contract give decoded output $F_{\rm XOP}(s_n,u_n)$,
which is the specified $s_{n+1}$. The restoration or reachable-representative
condition admits the next step. Induction proves agreement for the prefix.
\end{proof}
This proof identifies sufficient obligations. No gate layout, disturbance
envelope or memory implementation supplied with W1 establishes them yet.
It is therefore a conditional correctness result, not a hardware guarantee.

\subsection{A useful bit-margin lemma}
For one calibrated intensity decision, assume the noiseless low output is at
most $a_0$, the noiseless high output at least $a_1$, and additive readout error
is bounded in magnitude by $\epsilon\ge0$. Decode one when measured intensity
$I\ge\theta$. If
\begin{equation}
 a_0+\epsilon<\theta<a_1-\epsilon,
\end{equation}
then every admitted low is decoded zero and every admitted high is decoded one.
The proof is direct substitution of the worst allowed error into each bound.
The interval exists exactly when $a_1-a_0>2\epsilon$ for this strict-margin
profile. This is a tolerance derived from the declared gate model; it is not
a numerical tolerance already established for a device.

Restoring a digital bit under such a margin can preserve its exact logical
value while the underlying intensity varies continuously. It does not round
the mathematical operand represented by a whole exact expression stream.
The unresolved engineering question is whether all required gates, storage
and transfers maintain an adequate margin under the chosen physical conditions.
Distributions with unbounded tails need a probability statement instead of
a claim that a finite worst-case error bound holds on every trial.

\subsection{Finite-workload reliability without assumed independence}
For a finite workload with $N$ critical gate, memory or transfer events, let
$E_i$ be the event that the required local contract fails. Then
\begin{equation}
 \Pr\!\left(\bigcup_{i=1}^N E_i\right)
       \le\sum_{i=1}^N\Pr(E_i)\le\sum_{i=1}^N p_i.
\end{equation}
This union bound does not assume independence. For an adaptive computation,
define the disjoint first-failure events
$F_i=E_i\cap\bigcap_{j<i}E_j^c$ and pad inactive steps up to a declared finite
maximum $N$ with no-failure events. If every fault-free admitted history has
conditional probability of the next local failure at most $p_i$, and correct
steps preserve admission, then $\Pr(F_i)\le p_i$ by averaging over those
histories. Hence $\Pr(\text{any failure})=\sum_i\Pr(F_i)\le\sum_i p_i$.
No probability bound on already-failed, inadmissible histories is inferred.
Observed zero errors does not by itself establish $p_i=0$. A fixed nonzero
per-event estimate also does not prove reliability for arbitrarily long
execution. Source packet CRC and Boolean event agreement are separate checks
and cannot substitute for a hardware error model.

\subsection{What exact arithmetic does and does not remove}
The exact operator authority removes recurrent numerical rounding from the
specified expression evolution. It preserves the source's intentional
quantization when that quantization is part of the selected profile. It leaves
initial input error, physical model discrepancy, finite modal truncation,
calibration uncertainty, timing uncertainty, unresolved symbolic decisions and
physical execution faults as distinct quantities. These enter different
contracts and cannot all be assigned the name floating-point drift.

There is a concrete optimization strategy here: retain analytic geometry and
exact calibrated transfer expressions; share common graph nodes; eliminate a
square root only by a proved domain-correct identity; use geometric exclusion
only with the complete deformation bound; and consume each completed event
once. Final rounded exports can be produced when required without replacing
the expression roots. No claim of constant storage or constant computational
cost accompanies lossless one-bit encoding.
